%----------------------------------------------------------------------------------------
% RATIONALE VERGELIJKINGEN EN ONGELIJKHEDEN. Alle opgaven zijn eigen werk.
%
% Opbouw: Start, leerstof (theorie, werkwijze of uitleg, voorbeeld, korte oefening), en op een
% nieuwe pagina de Opdrachten (door elkaar, stijgende moeilijkheid), Koppel en Test jezelf.
% Labels van dit hoofdstuk beginnen met h3-.
%----------------------------------------------------------------------------------------

\mychapter{3}{Vergelijkingen en ongelijkheden}

\lkr{Leerplan B+S': LPD 18 en K1 (rationale functies), het deel over rationale vergelijkingen en ongelijkheden, ook in contexten. De twee klassieke fouten: een oplossing houden die de noemer nul maakt, en bij een ongelijkheid kruiselings vermenigvuldigen met een noemer waarvan het teken niet vastligt. Laat leerlingen beide fouten eens zelf maken en ontdekken (opdracht 7).}

Wanneer kost een stuk gemiddeld minder dan 4 euro? Hoe ver moet een voorwerp van een lens staan om een scherp beeld te geven? Zulke vragen geven vergelijkingen of ongelijkheden met de onbekende in een noemer. Je lost ze op met wat je over veeltermen en tekentabellen weet, als je op de noemers let.

\begin{oefeningen}[Start]
\oef{h3-start-1}

\oef{h3-verken}
\end{oefeningen}

%========================================================================================
\section{Rationale vergelijkingen}

Vermenigvuldig je beide leden met de noemer, dan verdwijnt de breuk. Maar dat mag alleen waar de noemer niet nul is: daarom schrijf je eerst de \textbf{bestaansvoorwaarden} op, en controleer je op het einde.

\voorbeeld{}{Los op: $\dfrac{1}{x} + \dfrac{1}{x - 1} = \dfrac32$.\label{vb:h3-vergelijking}}{%
\uitleg{Vermenigvuldig met het kleinste gemeenschappelijke veelvoud van de noemers, $2x(x - 1)$.}
\begin{stappenbreed}
x \neq 0,\ x \neq 1 & Bestaansvoorwaarden. \\
2(x - 1) + 2x = 3x(x - 1) & Vermenigvuldigen met $2x(x - 1)$. \\
3x^2 - 7x + 2 = 0 \iff x = 2 \text{ of } x = \tfrac13 & Discriminant $25$. \\
x \in \bigl\{\tfrac13, 2\bigr\} & Beide voldoen.
\end{stappenbreed}}

\voorbeeld{}{Los op: $\dfrac{x^2}{x - 2} = \dfrac{4}{x - 2}$.\label{vb:h3-vals}}{%
\uitleg{Na vermenigvuldigen vind je ook een getal dat de noemer nul maakt: dat is geen oplossing.}
\begin{stappenbreed}
x \neq 2 & Bestaansvoorwaarde. \\
x^2 = 4 \iff x = 2 \text{ of } x = -2 & Vermenigvuldigen met $x - 2$. \\
x = -2 & $x = 2$ voldoet niet.
\end{stappenbreed}}

\begin{werkwijze}
Schrijf de bestaansvoorwaarden op: geen enkele noemer mag nul zijn. Vermenigvuldig beide leden met het kgv van de noemers. Los de veeltermvergelijking op. Houd alleen de oplossingen die aan de bestaansvoorwaarden voldoen.
\end{werkwijze}

\begin{oefeningen}
\oef{h3-vergelijking}
\end{oefeningen}

%========================================================================================
\section{Rationale ongelijkheden}

Bij een ongelijkheid mag je niet zomaar met de noemer vermenigvuldigen: is die negatief, dan keert het teken om. Breng daarom alles naar één lid, maak er één breuk van, en gebruik een tekentabel.

\voorbeeld{}{Los op: $\dfrac{x + 1}{x - 2} \geqslant 3$.\label{vb:h3-ongelijkheid}}{%
\naast{%
\uitleg{Trek $3$ af, zet op één noemer en stel de tekentabel op. Een nulwaarde van de teller hoort bij de oplossing (want $\geqslant$), een nulwaarde van de noemer nooit.}
\begin{stappenvrij}{4.5cm}{3.2cm}
\frac{x + 1 - 3(x - 2)}{x - 2} \geqslant 0 & Eén lid, één breuk. \\[2.5ex]
\frac{-2x + 7}{x - 2} \geqslant 0 & Nulwaarde $\frac72$, pool $2$. \\[1.5ex]
x \in \left]2, \tfrac72\right] & Uit de tekentabel.
\end{stappenvrij}}{%
\tekentabel{cccccc}{x & & 2 & & \frac72 & \\\hline -2x + 7 & + & + & + & 0 & - \\ x - 2 & - & 0 & + & + & +}{\text{breuk} & - & | & + & 0 & -}}}

\begin{werkwijze}
Breng alle termen naar het linkerlid en zet ze op één noemer. Ontbind teller en noemer. Stel een tekentabel op met de nulwaarden van teller en noemer. Lees de oplossing af: de nulwaarden van de teller horen erbij bij $\leqslant$ of $\geqslant$, die van de noemer nooit.
\end{werkwijze}

\begin{oefeningen}
\oef{h3-ongelijkheid}
\end{oefeningen}

%========================================================================================
\section{Toepassingen}

\voorbeeld{}{Bij een productie van $q$ stuks is de gemiddelde kost $G(q) = 2 + \dfrac{500}{q}$ euro. Vanaf hoeveel stuks is ze lager dan $4$ euro?\label{vb:h3-context}}{%
\naast{%
\uitleg{Hier is $q > 0$, dus je mag wel met $q$ vermenigvuldigen zonder dat het teken omkeert.}
\begin{stappenvrij}{3.6cm}{3.6cm}
2 + \frac{500}{q} < 4 & \\[1ex]
500 < 2q & Maal $q > 0$. \\
q > 250 & Vanaf $251$ stuks.
\end{stappenvrij}}{%
\begin{tikzpicture}[x=.009cm,y=.18cm,font=\scriptsize]
  \draw[black!14,very thin,xstep=100,ystep=2] (0,0) grid (600,12);
  \draw[->,grijs] (0,0) -- (640,0) node[right,text=black] {$q$};
  \draw[->,grijs] (0,0) -- (0,13) node[above,text=black] {$G$};
  \foreach \t in {200,400,600}{\node[below,inner sep=1.5pt] at (\t,0) {$\t$};}
  \foreach \t in {2,4,8,12}{\node[left,inner sep=1.5pt] at (0,\t) {$\t$};}
  \draw[grijs,dashed] (0,2) -- (600,2);
  \draw[accent,dashed] (0,4) -- (600,4);
  \kromme{50}{600}{2+500/\x}
  \vol{(250,4)}
\end{tikzpicture}}}

\begin{oefeningen}
\oef{h3-context}
\end{oefeningen}

%========================================================================================
\opdrachtenkop

\begin{oefeningen}
\oef{h3-o1}

\oef{h3-o2}

\oef{h3-o3}

\oef{h3-o4}

\oef{h3-o5}

\oef{h3-o6}

\oef{h3-o7}

\oef{h3-o8}
\end{oefeningen}

%----------------------------------------------------------------------------------------
% KOPPEL: verbanden met andere leerstof
%----------------------------------------------------------------------------------------
\needspace{20\baselineskip}
\opdrachtsectie{Koppel}

\lkr{Kansrekenen en Rijen: kijk na wat al gezien is. Het pakket legt de volgorde niet vast.}

\begin{oefeningen}
\oef{h3-koppel-1}

\oef{h3-koppel-2}

\end{oefeningen}

%----------------------------------------------------------------------------------------
\opdrachtsectie{Test jezelf}

\begin{oefeningen}
\oef{h3-t1}

\oef{h3-t2}

\oef{h3-t3}

\oef{h3-t4}
\end{oefeningen}

\herhaalschema{%
  \fv{oef:h3-t1}{vb:h3-vergelijking}{oef:h3-vergelijking}%
  \fv{oef:h3-t2}{vb:h3-vals}{oef:h3-vergelijking}%
  \fv{oef:h3-t3}{vb:h3-ongelijkheid}{oef:h3-ongelijkheid}%
  \fv{oef:h3-t4}{vb:h3-context}{oef:h3-context}}

\ifzelfstudie
\section*{Maak je eigen samenvatting}
Sluit dit boek en schrijf in je schrift: de werkwijze voor rationale vergelijkingen met de bestaansvoorwaarden, waarom je bij een ongelijkheid niet zomaar met de noemer vermenigvuldigt, en de werkwijze met één breuk en een tekentabel. Vergelijk daarna met de kaders.
\fi
